\section{USAJMO 2012}
        \begin{enumerate}
        	\item (\textit{USAJMO 2012}) Given a triangle $ABC$, let $P$ and $Q$ be points on segments $\overline{AB}$ and $\overline{AC}$, respectively, such that $AP = AQ$.  Let $S$ and $R$ be distinct points on segment $\overline{BC}$ such that $S$ lies between $B$ and $R$, $\angle BPS = \angle PRS$, and $\angle CQR = \angle QSR$.  Prove that $P$, $Q$, $R$, $S$ are concyclic (in other words, these four points lie on a circle).


 



\textbf{Solution}

Since $\angle BPS = \angle PRS$, the circumcircle of triangle $PRS$ is tangent to $AB$ at $P$.  Similarly, since $\angle CQR = \angle QSR$, the circumcircle of triangle $QRS$ is tangent to $AC$ at $Q$.

 
                    \begin{figure}[H]
                        \centering
                        \includegraphics[width=0.5\textwidth]{Solution/USAJMO_I_2012_1_Solution_Figure_1.png}
                    \end{figure}
                    

 For the sake of contradiction, suppose that the circumcircles of triangles $PRS$ and $QRS$ are not the same circle.  Since $AP = AQ$, $AP^2 = AQ^2$, so the power of A with respect to both circumcircles is the same. Thus, $A$ lies on the \href{https://artofproblemsolving.com/wiki/index.php?title=Radical\_axis}{radical axis} of both circles.  However, both circles pass through $R$ and $S$, so the radical axis of both circles is $RS$.  Hence, $A$ lies on $RS$, which is a contradiction.

 Therefore, the two circumcircles are the same circle.  In other words, $P$, $Q$, $R$, and $S$ all lie on the same circle.

 



\textbf{Solution 2}

Note that (as in the first solution) the circumcircle of triangle $PRS$ is tangent to $AB$ at $P$.  Similarly, since $\angle CQR = \angle QSR$, the circumcircle of triangle $QRS$ is tangent to $AC$ at $Q$.

 Now, suppose these circumcircles are not the same circle. They already intersect at $R$ and $S$, so they cannot intersect anymore. Thus, AS must touch the two circumcircles at points $M$ and $N$, with $M$ on the circumcircle of triangle $PRS$. By Power of a Point, $AQ^2 = AM \cdot AS$ and $AP^2 = AN \cdot AS$. Hence, because $AP = AQ$, $AM = AN$, a contradiction because then, as they lie on the same line segment, M and N must be the same point! (Note line segment, not line.) Hence, the two circumcircles are the same circle.

 


	\item (\textit{USAJMO 2012}) Find all integers $n \ge 3$ such that among any $n$ positive real numbers $a_1$, $a_2$, $\dots$, $a_n$ with

 \[\max(a_1, a_2, \dots, a_n) \le n \cdot \min(a_1, a_2, \dots, a_n),\]
there exist three that are the side lengths of an acute triangle.


 



\textbf{Solution}

Without loss of generality, assume that the set $\{a\}$ is ordered from least to greatest so that the bounding condition becomes $a_n \le n \cdot a_1.$ Now set $b_i \equiv \frac{a_i}{a_1},$ and since a triangle with sidelengths from $\{a\}$ will be similar to the corresponding triangle from $\{b\},$ we simply have to show the existence of acute triangles in $\{b\}.$ Note that $b_1 = 1$ and for all $i$, $b_i \le n.$

 Now three arbitrary sidelengths $x$, $y$, and $z$, with $x \le y \le z,$ will form a valid triangle if and only if $x+y>z.$ Furthermore, this triangle will be acute if and only if $x^2 + y^2 > z^2.$ However, the first inequality can actually be inferred from the second, since $x+y>z \longleftrightarrow x^2 + y^2 +2xy > z^2$ and $2xy$ is trivially greater than $0.$ So we just need to find all $n$ such that there is necessarily a triplet of $b$'s for which $b_i^2 + b_j^2 > b_k^2$ (where $b_i < b_j < b_k$).

 We now make another substitution: $c_i \equiv b_i ^2.$ So $c_1 = 1$ and for all $i$, $c_i \le n^2.$ Now we examine the smallest possible sets $\{c\}$ for small $n$ for which the conditions of the problem are not met. Note that by smallest, we mean the set whose greatest element is as small as possible. If $n=3$, then the smallest possible set, call it $\{s_3\},$ is trivially $\{1,1,2\}$, since $c_1$ and $c_2$ are obviously minimized and $c_3$ follows as minimal. Using this as the base case, we see inductively that in general $\{s_n\}$ is the set of the first $n$ Fibonacci numbers. To show this note that if $\{s_n\} = \{F_0, F_1, ... F_n\}$, then $\{s_{n+1}\} = \{F_0, F_1, ... F_n, c_{n+1}\}.$ The smallest possible value for $c_{n+1}$ is the sum of the two greatest values of $\{s_n\}$ which are $F_{n-1}$ and $F_n$. But these sum to $F_{n+1}$ so $\{s_{n+1}\} = \{F_0, F_1, ... F_{n+1}\}$ and our induction is complete.

 Now since we know that the Fibonacci set is the smallest possible set which does not satisfy the conditions of the problem, then any set $\{c\}$ whose greatest term is less than $F_{n-1}$ must satisfy the conditions. And since $\{c\}$ is bounded between $1$ and $n^2$, then the conditions of the problem are met if and only if $F_{n-1} > n^2$. The first $n$ for which this restriction is satisfied is $n=13$ and the exponential behavior of the Fibonacci numbers ensure that every $n$ greater than $13$ will also satisfy this restriction. So the final solution set is $\boxed{\{n \ge 13\}}$.

 



\textbf{Solution 2}

\textbf{Outline:}

 1. Define the Fibonacci numbers to be $F_1 = F_2 = 1$ and $F_k = F_{k-1} + F_{k-2}$ for $k \ge 3$.

 2. If the chosen $n$ is such that $F_n \le n^2$, then choose the sequence $a_n$ such that $a_k = \sqrt{F_k}$ for $1 \le k \le n$. It is easy to verify that such a sequence satisfies the condition that the largest term is less than or equal to $n$ times the smallest term. Also, because for any three terms $x = \sqrt{F_a}, y = \sqrt{F_b}, z = \sqrt{F_c}$ with $a<b<c$, $x^2 + y^2 = F_a + F_b \le F_{b-1} + F_b = F_{b+1} \le F_c = z^2$, x, y, z do not form an acute triangle. Thus, all $n$ such that $F_n \le n^2$ do not work.

 3. It is easy to observe via a contradiction argument that all $n$ such that $F_n > n^2$ produce an acute triangle. (If, without loss of generality, $a_n$ is an increasing sequence, such that no three (in particular, consecutive) terms form an acute triangle, then $a_2^2 \ge F_1a_1^2, a_3^2 \ge a_2^2 + a_1^2 \ge F_2a^2$, and by induction $a_n^2 > F_na_1^2$, a contradiction to the condition's inequality.)

 4. Note that $F_{12} = 144 = 12^2$ and $F_{13} = 233 > 169 = 13^2$. It is easily to verify through strong induction that all $n$ greater than 12 make $F_n > n^2$. Thus, $\boxed{n \ge 13}$ is the desired solution set.

 


 


	\item (\textit{USAJMO 2012}) Let $a$, $b$, $c$ be positive real numbers.  Prove that

 \[\frac{a^3 + 3b^3}{5a + b} + \frac{b^3 + 3c^3}{5b + c} + \frac{c^3 + 3a^3}{5c + a} \ge \frac{2}{3} (a^2 + b^2 + c^2).\]


 



\textbf{Solution}

By the \href{https://artofproblemsolving.com/wiki/index.php?title=Cauchy-Schwarz\_Inequality}{Cauchy-Schwarz} inequality,
 \[[a(5a + b) + b(5b + c) + c(5c + a)] \left( \frac{a^3}{5a + b} + \frac{b^3}{5b + c} + \frac{c^3}{5c + a} \right) \ge (a^2 + b^2 + c^2)^2,\]so
 \[\frac{a^3}{5a + b} + \frac{b^3}{5b + c} + \frac{c^3}{5c + a} \ge \frac{(a^2 + b^2 + c^2)^2}{5a^2 + 5b^2 + 5c^2 + ab + ac + bc}.\]Since $a^2 + b^2 + c^2 \ge ab + ac + bc$,
 \[\frac{(a^2 + b^2 + c^2)^2}{5a^2 + 5b^2 + 5c^2 + ab + ac + bc} \ge \frac{(a^2 + b^2 + c^2)^2}{6a^2 + 6b^2 + 6c^2} = \frac{1}{6} (a^2 + b^2 + c^2).\]Hence,
 \[\frac{a^3}{5a + b} + \frac{b^3}{5b + c} + \frac{c^3}{5c + a} \ge \frac{1}{6} (a^2 + b^2 + c^2).\]

 Again by the Cauchy-Schwarz inequality,
 \[[b(5a + b) + c(5b + c) + a(5c + a)] \left( \frac{b^3}{5a + b} + \frac{c^3}{5b + c} + \frac{a^3}{5c + a} \right) \ge (a^2 + b^2 + c^2)^2,\]so
 \[\frac{b^3}{5a + b} + \frac{c^3}{5b + c} + \frac{a^3}{5c + a} \ge \frac{(a^2 + b^2 + c^2)^2}{a^2 + b^2 + c^2 + 5ab + 5ac + 5bc}.\]Since $a^2 + b^2 + c^2 \ge ab + ac + bc$,
 \[\frac{(a^2 + b^2 + c^2)^2}{a^2 + b^2 + c^2 + 5ab + 5ac + 5bc} \ge \frac{(a^2 + b^2 + c^2)^2}{6a^2 + 6b^2 + 6c^2} = \frac{1}{6} (a^2 + b^2 + c^2).\]Hence,
 \[\frac{b^3}{5a + b} + \frac{c^3}{5b + c} + \frac{a^3}{5c + a} \ge \frac{1}{6} (a^2 + b^2 + c^2).\]

 Therefore,
 \[\frac{a^3 + 3b^3}{5a + b} + \frac{b^3 + 3c^3}{5b + c} + \frac{c^3 + 3a^3}{5c + a} \ge \frac{1 + 3}{6} (a^2 + b^2 + c^2) = \frac{2}{3} (a^2 + b^2 + c^2).\]

 



\textbf{Solution 2}

Titu's Lemma: The sum of multiple fractions in the form $\frac{a_n^2}{b_n}$ where $a_n$ and $b_n$ are sequences of real numbers is greater than of equal to the square of the sum of all $a_i$ divided by the sum of all $b_i$, where i is a whole number less than n+1. Titu's Lemma can be proved using the Cauchy-Schwarz Inequality after multiplying out the denominator of the RHS.

 Consider the LHS of the proposed inequality. Split up the numerators and multiply both sides of the fraction by either a or 3a to make the LHS
 \[\sum_{cyc} \frac {a^4} {5a^2+ab}+\sum_{cyc} \frac {9a^4} {15ac+3a^2}\](Cyclic notation is a special notation in which all permutations of (a,b,c) is the summation command.)

 Then use Titu's Lemma on all terms: 
 \[\sum_{cyc} \frac {a^4} {5a^2+ab} +\sum_{cyc} \frac {9a^4} {15ac+3a^2} \ge \frac {16(a^2+b^2+c^2)^2} {8(a^2+b^2+c^2)+16(ab+bc+ca)} \ge \frac {16(a^2 + b^2 + c^2)^2}{24(a^2 + b^2 + c^2)} = \frac{2}{3} (a^2 + b^2 + c^2)\] owing to the fact that $a^2+b^2+c^2 \ge ab+bc+ca$, which is actually equivalent to $(a-b)^2 + (b-c)^2 + (c-a)^2 \ge 0$!

 



\textbf{Solution 3}

We proceed to prove that
 \[\frac{a^3 + 3b^3}{5a + b} \ge -\frac{1}{36} a^2 + \frac{25}{36} b^2\]

 (then the inequality in question is just the cyclic sum of both sides, since 
 \[\sum_{cyc} (-\frac{1}{36} a^2 + \frac{25}{36} b^2) = \frac{24}{36}\sum_{cyc} a^2 = \frac{2}{3} (a^2+b^2+c^2)\])

 Indeed, by AP-GP, we have 

 
 \[41 (a^3 + b^3+b^3) \ge 41 \cdot 3 ab^2\]

 and

 
 \[b^3 + a^2b \ge 2 ab^2\]

 Summing up, we have

 
 \[41a^3 + 83b^3 + a^2 b \ge 125 ab^2\]

 which is equivalent to:

 
 \[36(a^3 + 3b^3) \ge (5a + b)(-a^2 + 25b^2)\]

 Dividing $36(5a+b)$ from both sides, the desired inequality is proved.

 --\href{https://artofproblemsolving.com/wiki/index.php?title=User:Lightest}{Lightest} 15:31, 7 May 2012 (EDT)

 



\textbf{Solution 4}

By Cauchy-Schwarz, 

 
 \[\left(\sum_{cyc} \dfrac{a^3}{5a + b} + \dfrac{b^3}{5a + b} + \dfrac{b^3}{5a + b} + \dfrac{b^3}{5a + b} \right) \ge \dfrac{\left( \sum_{cyc} a^2 + b^2 + b^2 + b^2 \right)^2}{ \left( \sum_{cyc} a(5a + b) + b(5a + b) + b(5a + b) + b(5a + b) \right) }\]

 
 \[= \dfrac{\left(4a^2 + 4b^2 + 4c^2\right)^2}{\left(8a^2 + 8b^2 + 8c^2\right) + \left(16ab + 16bc + 16ca\right)}\]

 
 \[\ge \dfrac{16\left(a^2 + b^2 + c^2\right)^2}{\left(8a^2 + 8b^2 + 8c^2\right) + \left(8a^2 + 8b^2\right) + \left(8b^2 + 8c^2\right) + \left(8c^2 + 8a^2\right)}\] (by AM-GM) 
 \[= \dfrac{2}{3}\left(a^2 + b^2 + c^2\right)\] as desired.

 



\textbf{Solution 5}

We note that if we can prove that there exists a real number $x$ such that 
 \[\frac{a^3+3b^3}{5a+b}\geq xa^2+(\frac{2}{3}-x)b^2\] for all positive reals $a$ and $b$, then we are done since both sides are a cyclic sum of $a,b$, and $c$. Note that if we multiply both $a$ and $b$ by a constant $r$, the left and right sides will both multiply by $r^2$ (since the numerator on the left will multiply by $r^3$ and the denominator will multiply by $r$. This means that if the inequality is satisfied for an ordered pair $(a,b)$, then it is also satisfied for $(ar,br)$ for any positive real $r$. Thus, without loss of generality we can let $b=1$,and our inequality becomes 
 \[\frac{a^3+3}{5a+1}\geq a^2x+\frac{2}{3}-x.\] Expanding this and rearranging we get $a^3(1-5x)+a^2(-x)+a(5x-\frac{10}{3})+x+\frac{7}{3}\geq0$. Since $1$ is a root of the left side, we can factor this as 
 \[(a-1)(a^2(1-5x)+a(1-6x)-x-\frac{7}{3})\geq0.\] Note that for this to always be nonnegative, we must have $a^2(1-5x)+a(1-6x)-x-\frac{7}{3}$ be positive when $a>1$ and negative when $a<1$. Thus, it must have a root at $a=1$, so we plug in $a=1$ is a root and find that we must have $x=-\frac{1}{36}$. When $x=-\frac{1}{36}$, the expression factors as $(a-1)^2(\frac{41}{36}a+\frac{83}{36})\geq0$ which obviously is nonnegative for all positive reals $a$, so we are done.

 


 



\textbf{Solution 6}

$LHS=\sum_{cyc}\frac{a^3+3b^3}{5a+b}=\sum_{cyc}\frac{a^3}{5a+b}+\sum_{cyc}\frac{3b^3}{5a+b}=\sum_{cyc}\frac{a^4}{5a^2+ab}+\sum_{cyc}\frac{9b^4}{15ab+3b^2}$ which tells that $LHS\geq \frac{(a^2+b^2+c^2)^2}{5(a^2+b^2+c^2)+ab+bc+ac}+\frac{3(a^2+b^2+c^2)}{(a^2+b^2+c^2)+5(ab+bc+ac)}\geq (\frac{1}{2}+\frac{1}{6})(a^2+b^2+c^2)=RHS$, our proof is done

 


	\item (\textit{USAJMO 2012}) Let $\alpha$ be an irrational number with $0 < \alpha < 1$, and draw a circle in the plane whose circumference has length 1.  Given any integer $n \ge 3$, define a sequence of points $P_1$, $P_2$, $\dots$, $P_n$ as follows.  First select any point $P_1$ on the circle, and for $2 \le k \le n$ define $P_k$ as the point on the circle for which the length of arc $P_{k - 1} P_k$ is $\alpha$, when travelling counterclockwise around the circle from $P_{k - 1}$ to $P_k$.  Suppose that $P_a$ and $P_b$ are the nearest adjacent points on either side of $P_n$.  Prove that $a + b \le n$.


 



\textbf{Solution}

Use mathematical induction. For $n=3$ it is true because one point can't be closest to $P_3$ in both ways, and that $1+2\le 3$. Suppose that for some $n$, the nearest adjacent points $P_a$ and $P_b$ on either side of $P_n$ satisfy  $a+b \le n$. Then consider the nearest adjacent points $P_c$ and $P_d$ on either side of $P_{n+1}$. It is by the assumption of the nearness we can see that either $(c,d)=(a+1,b+1)$ still holds, or $P_1$ jumps into the interior of the arc $P_{a+1}P_{n}P_{b+1}$, so that $c$ or $d$ equals to $1$. Let's consider the following two cases.

 (i) Suppose $a+b=n$.

 Since the length of the arc $P_nP_a$ is $\{(a-n)\alpha\}$ (where $\{x\}$ equals to $x$ subtracted by the greatest integer not exceeding $x$) and length of the arc $P_bP_n$ is $\{(n-b)\alpha\} = \{a\alpha\}$, we now consider a point $P_0$ which is defined by $P_1$ traveling clockwise on the circle such that the length of arc $P_0P_1$ is $\alpha$. We claim that $P_0$ is in the interior of the arc $P_bP_nP_a$. Algebraically, it is equivalent to either $\{0-n\alpha\} < \{a\alpha-n\alpha\}$ or $\{n\alpha -0 \} < \{n\alpha - b\alpha\} = \{a\alpha\}$.

 Suppose the latter fails, i.e. $\{n\alpha\} \ge \{a\alpha\}$. Then suppose $n\alpha = m_1 + r_1$ and $a\alpha = m_2 + r_2$, where $m_1$, $m_2$ are integers and $0< r_2\le r_1 <1$ ($r_2$ is not zero because $a\alpha$ is irrational). We now have 
 \[\{0-n\alpha\} = \{-m_1-1 + (1 -r_1)\}=1-r_1\] and 
 \[\{a\alpha -n\alpha\} = \{m_2-m_1-1 + (1+r_2-r_1)\} = 1+r_2-r_1>1-r_1\]

 Therefore $P_0$ is either closer to $P_n$ than $P_a$ on the $P_a$ side, or closer to $P_n$ than $P_b$ on the $P_b$ side. In other words, $P_1$ is the closest adjacent point of $P_{n+1}$ on the $P_{a+1}$ side, or the closest adjacent point of $P_{n+1}$ on the $P_{b+1}$ side. Hence $P_c$ or $P_d$ is $P_1$, therefore $c+d \le n+1$.

 (ii) Suppose $a+b\le n-1$Then either $c+d = (a+1)+(b+1) \le n+1$ when $c=a+1$ and $d=b+1$, or $c+d \le 1+n$ when one of $P_c$ or $P_d$ is $P_1$.

 In either case, $c+d\le n+1$ is true.

 --\href{https://artofproblemsolving.com/wiki/index.php?title=User:Lightest}{Lightest} 20:27, 6 May 2012 (EDT)

 


	\item (\textit{USAJMO 2012}) For distinct positive integers $a$, $b < 2012$, define $f(a,b)$ to be the number of integers $k$ with $1 \le k < 2012$ such that the remainder when $ak$ divided by 2012 is greater than that of $bk$ divided by 2012.  Let $S$ be the minimum value of $f(a,b)$, where $a$ and $b$ range over all pairs of distinct positive integers less than 2012.  Determine $S$.


 



\textbf{Solution 1}

First we'll show that $S \geq 502$, then we'll find an example $(a, b)$ that have $f(a, b)=502$.

 Let $x_k$ be the remainder when $ak$ is divided by 2012, and let $y_k$ be defined similarly for $bk$. First, we know that, if $x_k > y_k >0$, then $x_{2012-k} \equiv a(2012-k) \equiv 2012-ak \equiv 2012-x_k \pmod {2012}$ and $y_{2012-k} \equiv 2012-y_k \pmod {2012}$. This implies that, since $2012 - x_k \neq 0$ and $2012 -y_k \neq 0$, $x_{2012-k} < y_{2012-k}$. Similarly, if $0< x_k < y_k$ then $x_{2012-k} > y_{2012-k}$, establishing a one-to-one correspondence between the number of $k$ such that $x_k < y_k$. Thus, if $n$ is the number of $k$ such that $x_k \neq y_k$ and $y_k \neq 0 \neq x_k$, then $S \geq \frac{1}{2}n$. Now I'll show that $n \geq 1004$.

 If $gcd(k, 2012)=1$, then I'll show you that $x_k \neq y_k$. This is actually pretty clear; assume that's not true and set up a congruence relation:
 \[ak \equiv bk \pmod {2012}\]Since $k$ is relatively prime to 2012, it is invertible mod 2012, so we must have $a \equiv b \pmod {2012}$. Since $0 < a, b <2012$, this means $a=b$, which the problem doesn't allow, thus contradiction, and $x_k \neq y_k$. Additionally, if $gcd(k, 2012)=1$, then $x_k \neq 0 \neq y_k$, then based on what we know about $n$ from the previous paragraph, $n$ is at least as large as the number of k relatively prime to 2012. Thus, $n \geq \phi(2012) = \phi(503*4) = 1004$. Thus, $S \geq 502$.

 To show 502 works, consider $(a, b)=(1006, 2)$. For all even $k$ we have $x_k=0$, so it doesn't count towards $f(1006, 2)$. Additionally, if $k = 503, 503*3$ then $x_k = y_k = 1006$, so the only number that count towards $f(1006, 2)$ are the odd numbers not divisible by 503. There are 1004 such numbers. However, for all such odd k not divisible by 503 (so numbers relatively prime to 2012), we have $x_k \neq 0 \neq y_k$ and $2012-k$ is also relatively prime to 2012. Since under those conditions exactly one of $x_k > y_k$ and $x_{2012-k} > y_{2012-k}$ is true, we have at most 1/2 of the 1004 possible k actually count to $f(1006, 2)$, so $\frac{1004}{2} = 502 \geq f(1006, 2) \geq S \geq 502$, so $S=502$.

 



\textbf{Solution 2}

Let $ak \equiv r_{a} \pmod{2012}$ and $bk \equiv r_{b} \pmod{2012}$. Notice that this means $a(2012 - k) \equiv 2012 - r_{a} \pmod{2012}$ and $b(2012 - k) \equiv 2012 - r_{b} \pmod{2012}$. Thus, for every value of $k$ where $r_{a} > r_{b}$, there is a value of $k$ where $r_{b} > r_{a}$. Therefore, we merely have to calculate $\frac{1}{2}$ times the number of values of $k$ for which $r_{a} \neq r_{b}$ and $r_{a} \neq 0$.

 
However, the answer is NOT $\frac{1}{2}(2012) = 1006$! This is because we must count the cases where the value of $k$ makes $r_{a} = r_{b}$ or where $r_{a} = 0$.

 
So, let's start counting.

 
If $k$ is even, we have either $a \equiv 0 \pmod{1006}$ or $a - b \equiv 0 \pmod{1006}$. So, $a = 1006$ or $a = b + 1006$. We have $1005$ even values of $k$ (which is all the possible even values of $k$, since the two above requirements don't put any bounds on $k$ at all).

 
If $k$ is odd, if $k = 503$ or $k = 503 \cdot 3$, then $a \equiv 0 \pmod{4}$ or $a \equiv b \pmod{4}$. Otherwise, $ak \equiv 0 \pmod{2012}$ or $ak \equiv bk \pmod{2012}$, which is impossible to satisfy, given the domain $a, b < 2012$. So, we have $2$ values of $k$.

 
In total, we have $2 + 1005 = 1007$ values of $k$ which makes $r_{a} = r_{b}$ or $r_{a} = 0$, so there are $2011 - 1007 = 1004$ values of $k$ for which $r_{a} \neq r_{b}$ and $r_{a} \neq 0$. Thus, by our reasoning above, our solution is $\frac{1}{2} \cdot 1004 = \boxed{502}$.

 
Solution by $\textbf{\underline{Invoker}}$

 



\textbf{Solution 3}

The key insight in this problem is noticing that when $ak$ is higher than $bk$, $a(2012-k)$ is lower than $b(2012-k)$, except at $2 \pmod{4}$ residues*. Also, they must be equal many times. $2012=2^2*503$. We should have multiples of $503$. After trying all three pairs and getting $503$ as our answer, we win. But look at the $2\pmod{4}$ idea. What if we just took $2$ and plugged it in with $1006$?We get $502$.

 <pre>--\href{https://artofproblemsolving.com/wiki/index.php?title=User:Va2010&action=edit&redlink=1}{Va2010} 11:12, 28 April 2012 (EDT)va2010</pre>



\textbf{Solution 4}

Say that the problem is a race track with $2012$ spots. To intersect the most, we should get next to each other a lot so the negation is high. As $2012=2^2*503$, we intersect at a lot of multiples of $503$.

 


	\item (\textit{USAJMO 2012}) Let $P$ be a point in the plane of triangle $ABC$, and $\gamma$ a line passing through $P$.  Let $A'$, $B'$, $C'$ be the points where the reflections of lines $PA$, $PB$, $PC$ with respect to $\gamma$ intersect lines $BC$, $AC$, $AB$, respectively.  Prove that $A'$, $B'$, $C'$ are collinear.


 



\textbf{Solution}

By the \href{https://artofproblemsolving.com/wiki/index.php?title=Law\_of\_Sines}{sine law} on triangle $AB'P$,
 \[\frac{AB'}{\sin \angle APB'} = \frac{AP}{\sin \angle AB'P},\]so
 \[AB' = AP \cdot \frac{\sin \angle APB'}{\sin \angle AB'P}.\]

 
                    \begin{figure}[H]
                        \centering
                        \includegraphics[width=0.5\textwidth]{Solution/USAJMO_I_2012_6_Solution_Figure_1.png}
                    \end{figure}
                    

 Similarly,
 \begin{align*} B'C &= CP \cdot \frac{\sin \angle CPB'}{\sin \angle CB'P}, \\ CA' &= CP \cdot \frac{\sin \angle CPA'}{\sin \angle CA'P}, \\ A'B &= BP \cdot \frac{\sin \angle BPA'}{\sin \angle BA'P}, \\ BC' &= BP \cdot \frac{\sin \angle BPC'}{\sin \angle BC'P}, \\ C'A &= AP \cdot \frac{\sin \angle APC'}{\sin \angle AC'P}. \end{align*}
Hence,
 \begin{align*} &\frac{AB'}{B'C} \cdot \frac{CA'}{A'B} \cdot \frac{BC'}{C'A} \\ &= \frac{\sin \angle APB'}{\sin \angle AB'P} \cdot \frac{\sin \angle CB'P}{\sin \angle CPB'} \cdot \frac{\sin \angle CPA'}{\sin \angle CA'P} \cdot \frac{\sin \angle BA'P}{\sin \angle BPA'} \cdot \frac{\sin \angle BPC'}{\sin \angle BC'P} \cdot \frac{\sin \angle AC'P}{\sin \angle APC'}. \end{align*}


 Since angles $\angle AB'P$ and $\angle CB'P$ are supplementary or equal, depending on the position of $B'$ on $AC$,
 \[\sin \angle AB'P = \sin \angle CB'P.\]Similarly,
 \begin{align*} \sin \angle CA'P &= \sin \angle BA'P, \\ \sin \angle BC'P &= \sin \angle AC'P. \end{align*}


 By the reflective property, $\angle APB'$ and $\angle BPA'$ are supplementary or equal, so
 \[\sin \angle APB' = \sin \angle BPA'.\]Similarly,
 \begin{align*} \sin \angle CPA' &= \sin \angle APC', \\ \sin \angle BPC' &= \sin \angle CPB'. \end{align*}
Therefore,
 \[\frac{AB'}{B'C} \cdot \frac{CA'}{A'B} \cdot \frac{BC'}{C'A} = 1,\]so by \href{https://artofproblemsolving.com/wiki/index.php?title=Menelaus\%27\_Theorem}{Menelaus's theorem}, $A'$, $B'$, and $C'$ are collinear.

 



\textbf{Solution 2, Barycentric (Modified by Evan Chen)}

We will perform barycentric coordinates on the triangle $PCC'$, with $P=(1,0,0)$, $C'=(0,1,0)$, and $C=(0,0,1)$. Set $a = CC'$, $b = CP$, $c = C'P$ as usual. Since $A$, $B$, $C'$ are collinear, we will define $A = (p : k : q)$ and $B = (p : \ell : q)$.

 Claim: Line $\gamma$ is the angle bisector of $\angle APA'$, $\angle BPB'$, and $\angle CPC'$. This is proved by observing that since $A'P$ is the reflection of $AP$ across $\gamma$, etc.

 Thus $B'$ is the intersection of the isogonal of $B$ with respect to $\angle P$with the line $CA$; that is,
 \[B' = \left( \frac pk \frac{b^2}{\ell}: \frac{b^2}{\ell} : \frac{c^2}{q} \right).\]Analogously, $A'$ is the intersection of the isogonal of $A$ with respect to $\angle P$with the line $CB$; that is,
 \[A' = \left( \frac{p}{\ell} \frac{b^2}{k} : \frac{b^2}{k} : \frac{c^2}{q} \right).\]The ratio of the first to third coordinate in these two pointsis both $b^2pq : c^2k\ell$, so it follows $A'$, $B'$, and $C'$ are collinear.

 



\textbf{Solution 3, Cartesian}

We will use Cartesian coordinates. Without loss of generality, let $P = (0,0)$ and let $\gamma$ be the line $x = 0$. Let $A = (x_1,y_1)$, $B = (x_2,y_2)$, and $C = (x_3,y_3)$. Then $A'$ is the intersection of the lines\begin{align*}y &= -\frac{y_1}{x_1}x \\y - y_2 &= \frac{y_2 - y_3}{x_2 - x_3}(x - x_2).\end{align*}
Solving the system of equations, we see it is
 \[\left(\frac{x_1(x_2y_3 - y_2x_3)}{x_1y_2 + y_1x_2 - x_1y_3 - x_3y_1},\frac{y_1(y_2x_3 - x_2y_3)}{x_1y_2 + y_1x_2 - x_1y_3 - x_3y_1}\right).\]To check collinearity, we need the determinant of the matrix
 \[\begin{bmatrix} x_1(x_2y_3 - y_2x_3) & y_1(y_2x_3 - x_2y_3) & x_1y_2 + y_1x_2 - x_1y_3 - x_3y_1 \\ x_2(x_3y_1 - y_3x_1) & y_2(y_3x_1 - x_3y_1) & x_2y_3 + y_2x_3 - x_2y_1 - x_1y_2 \\ x_3(x_1y_2 - y_1x_2) & y_3(y_1x_2 - x_1y_2) & x_3y_1 + y_3x_1 - x_3y_2 - x_2y_3 \end{bmatrix}\]to be zero, which is true because the columns sum to zero.

 


\end{enumerate}